% GENERATED FILE. Edit canonical lessons, then run npm run generate:books.
% canonical-source-hash: fnv1a64:c427a8e48fbf27fd
% canonical-lessons: JA-C82-machikado, JA-C82-shinai, JA-C82-kougai, JA-C82-hantaigawa, JA-C82-kawakami

\chapter{Street corner, Within the city, Suburbs, Opposite side, Upstream}
\label{ch:ja-street-corner-within-the-city-suburbs-opposite-side-upstream}
\hlchaptermodality{82}

\begin{chapteropening}
\textbf{By the end of this chapter:} \emph{I can say a street corner, within the city, the suburbs, the opposite side and upstream.}

Complete the last lesson of chapter 82: i can say a street corner, within the city, the suburbs, the opposite side and upstream.
\end{chapteropening}

\section[\texorpdfstring{\ja{まちかど} (machikado)}{まちかど (machikado)}]{\ja{まちかど} (machikado) — a street corner}

\label{lesson:JA-C82-machikado}



\begin{quote}
Before the new one: say the Japanese for between, then the Japanese for a region.
\end{quote}

\subsection*{You'll want to know: \ja{まちかど}}
\textbf{\ja{まちかど}} — \emph{machikado} — \textquotedblleft{}a street corner\textquotedblright{}.

\begin{grammarlens}[title={What you've built}]
One more word for this chapter.
\end{grammarlens}

\subsection*{Your turn}
\begin{itemize}
\raggedright
  \item \emph{Say it:} \emph{machikado}
  \item \emph{Say it:} \emph{machikado}, once more
  \item \emph{Say it:} say \emph{machikado}
  \item \emph{Recall:} say the Japanese for between, then the Japanese for a region, then say \emph{machikado} again
\end{itemize}

\begin{tcolorbox}[breakable,colback=teal!4,colframe=teal!35!black,title={Before you move on}]
What does \ja{まちかど} mean? (\textquotedblleft{}A street corner\textquotedblright{}.) Say it once more.
\end{tcolorbox}
\section[\texorpdfstring{\ja{しない} (shinai)}{しない (shinai)}]{\ja{しない} (shinai) — within the city}

\label{lesson:JA-C82-shinai}



\begin{quote}
Before the new one: say the Japanese for a region, then the Japanese for a street corner.
\end{quote}

\subsection*{You'll want to know: \ja{しない}}
\textbf{\ja{しない}} — \emph{shinai} — \textquotedblleft{}within the city\textquotedblright{}.

\begin{grammarlens}[title={What you've built}]
One more word for this chapter.
\end{grammarlens}

\subsection*{Your turn}
\begin{itemize}
\raggedright
  \item \emph{Say it:} \emph{shinai}
  \item \emph{Say it:} \emph{shinai}, once more
  \item \emph{Say it:} say \emph{shinai}
  \item \emph{Recall:} say the Japanese for a region, then the Japanese for a street corner, then say \emph{shinai} again
\end{itemize}

\begin{tcolorbox}[breakable,colback=teal!4,colframe=teal!35!black,title={Before you move on}]
What does \ja{しない} mean? (\textquotedblleft{}Within the city\textquotedblright{}.) Say it once more.
\end{tcolorbox}
\section[\texorpdfstring{\ja{こうがい} (kougai)}{こうがい (kougai)}]{\ja{こうがい} (kougai) — the suburbs}

\label{lesson:JA-C82-kougai}



\begin{quote}
Before the new one: say the Japanese for a street corner, then the Japanese for within the city.
\end{quote}

\subsection*{You'll want to know: \ja{こうがい}}
\textbf{\ja{こうがい}} — \emph{kougai} — \textquotedblleft{}the suburbs\textquotedblright{}.

\begin{grammarlens}[title={What you've built}]
One more word for this chapter.
\end{grammarlens}

\subsection*{Your turn}
\begin{itemize}
\raggedright
  \item \emph{Say it:} \emph{kougai}
  \item \emph{Say it:} \emph{kougai}, once more
  \item \emph{Say it:} say \emph{kougai}
  \item \emph{Recall:} say the Japanese for a street corner, then the Japanese for within the city, then say \emph{kougai} again
\end{itemize}

\begin{tcolorbox}[breakable,colback=teal!4,colframe=teal!35!black,title={Before you move on}]
What does \ja{こうがい} mean? (\textquotedblleft{}The suburbs\textquotedblright{}.) Say it once more.
\end{tcolorbox}
\section[\texorpdfstring{\ja{はんたいがわ} (hantaigawa)}{はんたいがわ (hantaigawa)}]{\ja{はんたいがわ} (hantaigawa) — the opposite side}

\label{lesson:JA-C82-hantaigawa}



\begin{quote}
Before the new one: say the Japanese for within the city, then the Japanese for the suburbs.
\end{quote}

\subsection*{You'll want to know: \ja{はんたいがわ}}
\textbf{\ja{はんたいがわ}} — \emph{hantaigawa} — \textquotedblleft{}the opposite side\textquotedblright{}.

\begin{grammarlens}[title={What you've built}]
One more word for this chapter.
\end{grammarlens}

\subsection*{Your turn}
\begin{itemize}
\raggedright
  \item \emph{Say it:} \emph{hantaigawa}
  \item \emph{Say it:} \emph{hantaigawa}, once more
  \item \emph{Say it:} say \emph{hantaigawa}
  \item \emph{Recall:} say the Japanese for within the city, then the Japanese for the suburbs, then say \emph{hantaigawa} again
\end{itemize}

\begin{tcolorbox}[breakable,colback=teal!4,colframe=teal!35!black,title={Before you move on}]
What does \ja{はんたいがわ} mean? (\textquotedblleft{}The opposite side\textquotedblright{}.) Say it once more.
\end{tcolorbox}
\section[\texorpdfstring{\ja{かわかみ} (kawakami)}{かわかみ (kawakami)}]{\ja{かわかみ} (kawakami) — upstream}

\label{lesson:JA-C82-kawakami}



\begin{quote}
Before the new one: say the Japanese for the suburbs, then the Japanese for the opposite side.
\end{quote}

\subsection*{You'll want to know: \ja{かわかみ}}
\textbf{\ja{かわかみ}} — \emph{kawakami} — \textquotedblleft{}upstream\textquotedblright{}.

\begin{grammarlens}[title={What you've built}]
One more word for this chapter.
\end{grammarlens}

\subsection*{Your turn}
\begin{itemize}
\raggedright
  \item \emph{Say it:} \emph{kawakami}
  \item \emph{Say it:} \emph{kawakami}, once more
  \item \emph{Say it:} say \emph{kawakami}
  \item \emph{Recall:} say the Japanese for the suburbs, then the Japanese for the opposite side, then say \emph{kawakami} again
\end{itemize}

\begin{tcolorbox}[breakable,colback=teal!4,colframe=teal!35!black,title={Before you move on}]
What does \ja{かわかみ} mean? (\textquotedblleft{}Upstream\textquotedblright{}.) Say it once more.
\end{tcolorbox}
